\documentclass[11pt]{article}
\usepackage[a4paper, top=18mm, bottom=18mm, left=20mm, right=20mm]{geometry}
\usepackage[T2A]{fontenc}
\usepackage[utf8]{inputenc}
\usepackage[ukrainian]{babel}
\usepackage{graphicx}
\usepackage[export]{adjustbox}
\usepackage{amsmath}
\usepackage{amssymb}
\setlength{\parindent}{0pt}
\setlength{\parskip}{4pt}
\pagestyle{empty}
\begin{document}
%begin
{\large\textbf{Контрольна робота}}

\textbf{Варіант \No\,3}\hfill \textit{ПІБ:}\,\underline{\hspace{60mm}}
\vspace{4mm}

%beginitem
1. Що буде виведено за виконання фрагменту коду?

%code
for f in range(2,2+4,3):
    if f>=8:
        pass
    print(f, end=' ')
    f=2
    if f>=8:
        break
else:
    print(f, end=' ')
print(f, end=' ')
%code

%enditem

%beginitem
2. 
Змінні \kw{next} та \kw{prev} вузла списку позначають наступний та попередній за ним вузол відповідно. Починаючи з голови, у список записано послідовні цілі числа від~$-21$ до~$20$. Нехай змінна \kw{p} позначає вузол, що зберігає значення~$3$.
Яке число зберігається у вузлі \kw{p.next.next.next.next.next} ?

%enditem

%beginitem
3. 
Записати ВИРАЗ, значенням якого є множина, що складається з кубівцілих чисел від 12 до 2012 (включно), що не діляться на жодне з чисел 2, 3.
%enditem

%beginitem
4. 
try:
    t = 5
    while t<=7:
        t += 2
    print(t, end=';')
except:
    print('ERR')

%enditem

%beginitem
5. 
def f(n): print("f", end=""); return n!=-2
   print(f(3) or f(-6))
%enditem

%beginitem
6. %goodpath:Scripts/2019/Mod2019_1/Lex/03-good.txt
Відмітити все, що є коректним оператором С++:
( 1 ) x not is y

( 2 ) print(end="1","qwe")

( 3 ) 3**-1

( 4 ) x-=2

%enditem

%beginitem
7. %goodpath:Scripts/2019/Mod2019_1/Lex/01-good.txt
Відмітити все, що є однією правильною лексемою:
( 1 ) 5_w

( 2 ) 0x17

( 3 ) 'c'

( 4 ) "1,3,2"

%enditem

%beginitem
8. 

def h(a,b):
    c=84
    if b!=-3:
        c=5
    elif a>-2:
         return 2
    else: 
        return 0
    return c

print(h(3,3))
%enditem

%beginitem
9. %goodpath:Scripts/2018/Mod2018_1/01-good.txt
Відмітити все, що є однією правильною лексемою:
( 1 ) 5_w

( 2 ) (1)

( 3 ) )

( 4 ) 0.5

%enditem

%beginitem
10. 
for f in range(-3,-3,-3):
    if f>=8:
        break
        print(f,end=' ')
    if f>=8:
        break
else:
    print(f, end=' ')
print(f, end=' ')

%enditem

%end
\newpage
\end{document}

